\section{Thiết kế, dữ liệu và mức bằng chứng}


\subsection{Đầu vào và đối tượng so sánh}

Nhu cầu được lấy từ các chuyến taxi công khai của New York TLC~\cite{tlc}, qua bộ dữ
liệu Manhattan đóng gói cho FleetPy~\cite{fleetpy2022}. Mỗi ô là một tổ hợp ngày, mẫu nhu
cầu và khung giờ. Panel chính gồm năm ngày 14--18/11/2018, mỗi ngày bốn mẫu, khung
08:00--10:00. Các ngày này gồm ba ngày làm việc và hai ngày cuối tuần. Mỗi ô có 108 yêu
cầu, 8 xe bốn chỗ và cửa sổ đón 600 giây; tổng 20 ô có 2.160 yêu cầu.

Mạng thử không phải mạng phố Manhattan đầy đủ. Công cụ chuẩn hóa chọn 96 điểm và dựng
bao đóng thời gian đi, với một cung cho mỗi cặp điểm khác nhau: $96\times95=9.120$ cung.
Tọa độ nút là tổng hợp. Đây là một thực thể nhỏ để kiểm cơ chế trên cùng đầu vào; khóa luận
không gọi nó là thí nghiệm vận hành quy mô Manhattan.

Trong panel chính, thời gian đi được giữ bằng một ảnh chụp. Xe vẫn có tiến độ và thời gian
dừng trong mô phỏng rời rạc, nên phần thay đổi khi giữ tuyến không nhất thiết bằng không.
Các đợt có trôi dạt dùng hệ số thời gian đi theo ngày, mỗi 900 giây, áp chung cho mạng.
Chúng được tách khỏi panel tĩnh vì mô hình thông tin và các quy tắc thực thi khác nhau.

\begin{table}[htbp]
\centering
\caption{Các tập bằng chứng dùng để trả lời câu hỏi nghiên cứu}
\label{bang:lotrinh}
\small
\begin{tabularx}{\linewidth}{@{}p{2.5cm}p{3.4cm}Xp{3.2cm}@{}}
\toprule
\textbf{Tập} & \textbf{Đơn vị và phạm vi} & \textbf{Câu hỏi} & \textbf{Mức bằng chứng}\\
\midrule
H6 chính & 20 ô tĩnh, 2.160 yêu cầu mỗi nhánh & Ba mốc và chi phí của gói cam kết & Phép so chính đã niêm phong\\
Thang ngân sách & Cùng 20 ô; cửa sổ 600 và 900 s & Giá theo trần và cửa sổ & Phân tích bổ sung\\
Tầng 2 & 20 ô chính; 16 ô phát triển sáng/tối & Bốn họ cổng và nguyên nhân dừng & Thăm dò, dự đoán niêm phong trước\\
W07/W07X & 8 ô 07:00--09:00, ngày 12--13/11 & Lời hứa chặt và phục hồi & Phát triển, thăm dò có niêm phong\\
T38 & 15 ô hiệu chỉnh, 34 ô kiểm tra & Số khách, miễn trừ và biên nhìn thấy & Thăm dò, ngưỡng cố định trước tập kiểm tra\\
\bottomrule
\end{tabularx}
\end{table}

Bảng~\ref{bang:lotrinh} chỉ đếm các đơn vị cần cho lập luận, không lấy tổng số lần chạy hay
dung lượng lưu trữ làm mức mạnh của bằng chứng. Chạy nhiều cấu hình trên cùng một ô không
biến ô đó thành nhiều mẫu độc lập. Hạt giống của bộ điều phối được giữ bằng 7 trong các
phép so; lặp cùng cấu hình kiểm tính tái lập, không ước lượng biến thiên giữa các hạt giống.

\subsection{Cấu hình và phép ghép cặp}

Ở panel chính, \textbf{B1} không bật sổ ngân sách cam kết; nó vẫn tuân theo các ràng buộc
vật lý và không cho tái gán khách đã nhận sang xe khác. Vì vậy B1 không phải đối chứng
``tái gán tự do''. \textbf{C1} bật gói khóa xác nhận và trần hành vi 30 giây trên tổng độ
dời giờ trả. Các nhánh ngân sách khác giữ gói khóa và chỉ thay trần giờ trả. Nhánh
\texttt{nobudget} bỏ trần cộng dồn, nhưng \emph{vẫn khóa giờ đón}, cùng xe và điểm đón.

So B1 với C1 đo chi phí của cả gói, không riêng trần 30 giây. So \texttt{nobudget} với các
mức ngân sách gần hơn với câu hỏi về trần giờ trả, vì các khóa xác nhận được giữ. Cùng ô
dùng cùng yêu cầu, đội xe, bộ sinh và hạt giống. Dù vậy, các nhánh có thể chọn tuyến khác
nhau, kéo theo trạng thái và cơ hội ghép về sau khác nhau.

Ở tầng thực nghiệm mới, $U$ là nhánh không cổng thời gian đang xét; $C$-$x$, $M^{+}$-$x$
và $V$-$x$ lần lượt bật trần hành vi, hạn chót từ lời hứa đầu và trần nhìn thấy. Các khóa
và điều kiện vật lý khác được giữ theo cấu hình của đợt. $U$ ở T38 không được ghép trực
tiếp với B1 của H6, vì runner, thế giới và quy tắc ghi nhận vi phạm đón khác nhau.

\subsection{Chỉ số và đối chiếu nguồn}

Tỉ lệ phục vụ là số yêu cầu hoàn tất chia toàn bộ yêu cầu của ô. Chênh tỉ lệ có cùng mẫu số
được báo bằng điểm phần trăm. Biến động được tính từ chuỗi lời hứa, với tổng $C$, $X$, $V$
định nghĩa ở Chương~\ref{ch:phuongphap}. Với H6, phép đo trung tâm cộng cả hai tọa độ đón
và trả, chỉ khi tọa độ tương ứng còn mở. Các phép thử cổng và T38 tập trung vào giờ trả;
không trộn hai phạm vi thành một chỉ số.

Tỉ lệ thất hứa của một luật phải dùng \emph{đại lượng của luật đó}: $C_r>x$ cho cổng hành
vi; giờ trả thật muộn hơn lời hứa đầu quá $x$ cho hạn chót; $V_r>x$ cho cổng nhìn thấy.
Tỉ lệ này không phải một thang trải nghiệm chung giữa ba luật. Chỉ số chung E2 là tỉ lệ
khách đã được hứa có $V_r>60$ giây. Ngưỡng là dấu $>$ nghiêm ngặt, không phải $\ge$.

Một tỉ lệ có thể được gộp theo khách hoặc lấy trung bình không trọng số theo ô. Hai cách
có thể cho số khác nhau khi số khách mỗi ô khác nhau. Khóa luận ghi cách gộp cạnh mỗi bảng.
Tỉ lệ quyết định có miễn trừ dùng số quyết định đã được xác nhận làm mẫu số; không so
trực tiếp nó với tỉ lệ khách vượt lời hứa.

Trong lần biên tập này, phép đo H6 được đối chiếu lại với 40 biên bản gốc của B1/C1 và
điểm kiểm tra cuối của 180 lần chạy thuộc chín nhánh. Số phục vụ trên hai cửa sổ được
tính lại từ kết quả từng ô. Với T38, cả 185 điểm kiểm tra cuối được giải mã, kiểm băm
khung và tái tính lịch sử. Mục~\ref{muc:t38} báo phạm vi kiểm delta riêng. Những kiểm
tra này xác nhận số liệu mô tả; chúng không biến kết quả thăm dò thành xác nhận độc lập.

\section{RQ1: đại lượng được ghi sổ có phản ánh quyết định?}
\label{muc:quytrach}

\subsection{Tỉ số trung tâm và cách đọc}

Câu hỏi ở đây là phần khác biệt giữa phương án mới với tuyến giữ nguyên lớn bao nhiêu so
với biến động khách nhìn thấy. Bảng~\ref{bang:quytrach} tổng hợp theo khách đã được nhận,
cộng hai tọa độ đón và trả khi còn mở, trên toàn 20 ô mỗi nhánh. Đơn vị của các tổng là
giây; giá trị bảng đã làm tròn từ mili giây nguyên.

\begin{table}[htbp]
\centering
\caption{Các tổng biến động của B1 và C1 trên panel H6}
\label{bang:quytrach}
\begin{tabularx}{\linewidth}{@{}Xrr@{}}
\toprule
\textbf{Đại lượng} & \textbf{B1} & \textbf{C1}\\
\midrule
Yêu cầu hoàn tất & 1.735 & 1.581\\
Tỉ lệ phục vụ (\%) & 80,32 & 73,19\\
Số lần công bố lời hứa & 54.311 & 45.542\\
Tổng biến động nhìn thấy (s) & 83.587,440 & 9.332,101\\
Tổng phần do quyết định (s) & 74.443,002 & 128,020\\
Tổng phần thay đổi khi giữ tuyến (s) & 9.202,610 & 9.204,481\\
Tỉ số tổng do quyết định / tổng nhìn thấy (\%) & 89,06 & 1,37\\
Khách có phần do quyết định khác không & 403 & 11\\
Tổng độ dời hành vi lớn nhất theo khách (s) & 840,910 & 27,588\\
\bottomrule
\end{tabularx}
\end{table}

Tỉ số 89,06\% được tính bằng $74.443{,}002/83.587{,}440$. Nó cho thấy trong thiết lập
B1, tổng phần do quyết định gần bằng chín phần mười tổng biến động nhìn thấy. Đây là tỉ
số giữa hai tổng, \emph{không} phải 89,06\% tác động nhân quả của chính sách, và không
bảo đảm mỗi khách đều có tỉ số đó. Do có triệt tiêu dấu, tổng do quyết định cộng tổng giữ
tuyến không đúng bằng tổng nhìn thấy.

Dưới C1, tổng phần do quyết định còn 128,020 giây, trong khi phần giữ tuyến vẫn khoảng
9.204 giây. Tỉ số xuống 1,37\% cho thấy trần đã làm biến động còn lại chủ yếu nằm ngoài
phần bị tính phí. Nó cũng chỉ ra giới hạn thực dụng: giảm mạnh sổ hành vi không làm giờ
hẹn trên ứng dụng hoàn toàn bất động.

Hai tổng giữ tuyến của B1/C1 gần bằng nhau không tạo thành một ``nhóm chứng tự nhiên''
chứng minh nhân quả. Tuyến đã chọn có thể tác động tiến độ xe và số lần công bố. Trên
những cấu hình khác, tổng này cũng thay đổi. Điểm được kiểm ở đây là cách hạch toán và
phản ứng của các đại lượng trong panel, không phải tính độc lập của một nhiễu ngoại sinh.

\subsection{Biến động tập trung và tính không cộng của trị tuyệt đối}

Ở B1, 403/1.735 khách, tức 23,2\%, có phần do quyết định khác không. Khách có tổng độ
dời hai tọa độ lớn nhất đạt 840,910 giây. Đây là gần 14 phút \emph{tổng dịch chuyển lời
hứa}, không phải 14 phút trễ đến nơi hay chờ thêm: cộng hai tọa độ và nhiều lần sửa có
thể tính cùng một thay đổi ở nhiều thời điểm. Dưới C1, chỉ 11/1.581 khách có độ dời hành
vi khác không. Trần theo từng khách bảo vệ trường hợp lớn nhất mà một trung bình toàn
đội có thể làm mờ.

Kiểm riêng tọa độ giờ trả làm rõ vì sao cần bất đẳng thức. Ở B1, trong 54.311 lần công
bố, 228 lần có $A+d>B$; ở C1 có 3/45.542 lần. Trong 476 lần của B1 mà cả $A$ và $d$
đều khác không, dấu bằng hỏng ở 228 lần, tức 47,9\%. Vì vậy tỉ lệ dấu bằng đúng trên
mọi lần công bố rất cao phần lớn do một thành phần bằng không, không phải một kiểm
chứng mạnh cho đẳng thức phân rã.

Những phép kiểm dấu bằng không kiểm trực tiếp giả thiết (H), $A\le B$, của Định
lý~\ref{dl:tapcon}; dấu bằng hỏng chưa tự động có nghĩa $A>B$. Hơn nữa, bản ghi chỉ
chứa phương án đã chọn, trong khi (H) cần đúng trên mọi quỹ đạo được so. Không được
coi tỉ lệ dấu bằng cao là bằng chứng đã thỏa giả thiết đó. Quan hệ bao hàm có điều
kiện chỉ là công cụ hiểu mô hình; so số khách của các nhánh vẫn cần dữ liệu riêng.

\subsection{Phát biểu được hỗ trợ}

RQ1 có hai câu trả lời tách biệt. Về cài đặt, các mốc trước, giữ tuyến và đề xuất cho phép
tái tính phần được ghi sổ ở thời điểm quyết định. Về quan sát, gói C1 giảm mạnh phần này
trên cùng panel. Không có dữ liệu hài lòng, hành vi hủy hoặc lựa chọn người dùng để kết
luận rằng mức giảm đó cải thiện trải nghiệm một lượng xác định. Vì thế ``giảm biến động
hành vi'' là kết quả kỹ thuật, còn lợi ích người dùng là câu hỏi chưa được đo.

\section{RQ2: giá dịch vụ của gói cam kết và của trần}


\subsection{Gói chính và tiêu chí H6}

H6 là phép so chính được niêm phong, với yêu cầu đồng thời về giảm gánh sửa lời hứa và
giữ chi phí dịch vụ trong mức đã đặt. B1 hoàn tất 1.735/2.160 yêu cầu; C1 hoàn tất
1.581/2.160. Chênh 154 yêu cầu tương ứng $154/2.160\times100=7{,}13$ điểm phần trăm.
Đây là chi phí có thật, không phải phần trăm giảm tương đối so với B1.

Tiêu chí dịch vụ của H6 chỉ cho phép mất tối đa 1 điểm phần trăm. C1 không đạt tiêu chí
đó, dù phần kiểm ngân sách và giảm gánh sửa lời hứa đạt. Vì giả thuyết chính cần cả hai,
khóa luận ghi nhận \emph{kết quả âm đối với tiêu chí tổng hợp H6}. Việc bổ sung kết quả
thăm dò về sau không làm tiêu chí cũ trở thành đạt. Các đợt chỉnh adapter và chạy lại đã
giữ các số phục vụ trung tâm; không có bằng chứng chi phí 7,13 điểm đã biến mất.

Không thể quy 7,13 điểm riêng cho trần giờ trả 30 giây. Gói C1 còn gồm khóa xác nhận.
Trong các cấu hình được đo, nhánh không trần cộng dồn nhưng còn khóa xác nhận hoàn tất
1.655 yêu cầu, mất 3,70 điểm so B1. Vì nhánh này vẫn khóa giờ đón, kết quả không chứng
minh ``chỉ khóa danh tính đã tốn hơn nửa chi phí trước khi hứa về thời gian''.

\subsection{Thang ngân sách với các khóa được giữ}

Bảng~\ref{bang:thang1} giữ số chính xác để đọc giá ở mỗi mức. Phần do quyết định cộng
cả giờ đón và trả; với các khóa giờ đón của thang này, biến động chủ yếu nằm ở giờ trả.
Cột trần không áp cho B1; không được đọc giá trị lớn nhất của B1 như một vi phạm ngân
sách chưa được cấu hình.

\begin{table}[htbp]
\centering
\caption{Thang ngân sách hành vi, cửa sổ đón 600 giây, 20 ô}
\label{bang:thang1}
\small
\begin{tabular}{@{}lrrrr@{}}
\toprule
\textbf{Cấu hình} & \textbf{Hoàn tất} & \textbf{Mất so B1 (pp)} & \textbf{Tổng hành vi (s)} & \textbf{Lớn nhất (s)}\\
\midrule
B1 & 1.735 & 0,00 & 74.443,0 & 840,9\\
Khóa, không trần & 1.655 & 3,70 & 22.485,2 & 595,3\\
300 s & 1.641 & 4,35 & 15.808,8 & 290,8\\
120 s & 1.608 & 5,88 & 5.063,3 & 119,0\\
60 s & 1.579 & 7,22 & 1.108,8 & 59,0\\
30 s (C1) & 1.581 & 7,13 & 128,0 & 27,6\\
\bottomrule
\end{tabular}
\end{table}

Các mức có trần đều giữ giá trị lớn nhất theo khách dưới trần: 290,845/300;
118,958/120; 58,981/60 và 27,588/30 giây. Đây là đối chiếu trực tiếp của sổ cái trong
những lần chạy bình thường được xét. Từ không trần xuống 300 giây, số hoàn tất giảm
14; từ 300 xuống 120 giảm 33; từ 120 xuống 60 giảm 29. Mỗi bước mua thêm giới hạn cho
khách chịu biến động lớn nhất, nhưng giá dịch vụ không tuyến tính theo số giây của trần.

Trần 30 giây hoàn tất nhiều hơn 60 giây đúng hai khách. So từng ô, có khác biệt theo cả
hai chiều; tám trong 20 ô không theo thứ tự đơn điệu của mức trần. Đây là độ nhạy quỹ
đạo của bộ chọn cận thị trên đầu vào cố định. Không có cơ sở đặt một ``sàn nhiễu''
$\pm0{,}3$ điểm phần trăm rồi gọi những bước dưới sàn là thuần nhiễu. Chênh hai khách
nhỏ và không hỗ trợ một ưu thế dịch vụ có thể khái quát cho trần 30 giây.

\subsection{Cửa sổ rộng hơn có làm cam kết rẻ hơn?}



Phép kiểm bổ sung nới cửa sổ đón từ 600 lên 900 giây, giữ cùng panel. Cửa sổ rộng hơn
làm B1 hoàn tất 1.938 yêu cầu thay vì 1.735; C1 đạt 1.753 thay vì 1.581. Dịch vụ tăng
ở cả hai, nhưng khoảng cách tăng từ 154 lên 185 yêu cầu, tức từ 7,13 lên 8,56 điểm phần
trăm. Nới cửa sổ không làm giá của gói cam kết co lại trong phép kiểm này.

Hình~\ref{hinh:ngansach} so toàn thang trần bằng mức mất so với B1 \emph{của từng cửa
sổ}. Cách chuẩn hóa này tránh lẫn tăng dịch vụ do cửa sổ với tăng chi phí tương đối của
gói. Ở cả năm mức, khoảng cách so B1 của cửa sổ 900 giây lớn hơn của cửa sổ 600 giây.
Chênh khoảng cách ở trần 300 giây là 0,42 điểm; ở trần 30 giây là 1,44 điểm.

\begin{figure}[htbp]
\centering
\includegraphics[width=\linewidth]{hinh/service_budget_final.pdf}
\caption{Chi phí dịch vụ theo trần hành vi, so với B1 của từng cửa sổ đón. Mỗi điểm gộp
2.160 yêu cầu của cùng 20 ô; đường chỉ nối các cấu hình, không là hàm khớp hay khoảng bất
định. Điểm không trần vẫn giữ gói khóa xác nhận.}
\label{hinh:ngansach}
\end{figure}

Quan sát này phù hợp với khả năng cửa sổ rộng tạo thêm lựa chọn mà trần chặt không cho
khai thác. Tuy nhiên, hai mức cửa sổ chưa chứng minh một quy luật rằng mọi loại tự do
đội xe đều làm lời hứa đắt hơn. Tăng số xe, đổi mật độ nhu cầu hoặc cho tái gán có thể
tác động khác. Kết quả được giữ như độ nhạy theo cửa sổ, không nâng thành định luật
chung của các hệ điều phối.

Nới cửa sổ cũng thay chất lượng dịch vụ của người được nhận. Ở cấu hình C1, thời gian
chờ trung vị tăng từ 407 lên 636 giây trong phân tích bổ sung. Vì vậy tăng số khách
hoàn tất không đủ để gọi cấu hình 900 giây tốt hơn. Khóa luận không quy tổng độ dời
lời hứa thành tiền hay phúc lợi, vì chưa có hệ số định giá đúng đại lượng và quần thể
đang xét.

\subsection{Cách trả lời RQ2}

Trên panel chính, gói C1 mua được một trần hành vi kiểm được và giảm mạnh tổng biến
động do quyết định, với chi phí 7,13 điểm phần trăm dịch vụ. Đây là một đánh đổi, đồng
thời là thất bại đối với mức chi phí H6 đã đặt. Thang ngân sách giúp chọn một điểm
vận hành theo ưu tiên cụ thể; không có một mức trần tốt nhất cho mọi nhu cầu.

Giá của trần phải được tách khỏi giá của các khóa đi kèm. Việc chọn trần cũng cần xét
cửa sổ dịch vụ và biến động nhìn thấy. Chỉ đặt một trần nhỏ rồi báo số chuyến giảm
không đủ để kết luận khách sẽ hài lòng hơn hoặc người vận hành nên chấp nhận chi phí
đó. Những quyết định này cần dữ liệu người dùng và điều kiện vận hành chưa có trong
khóa luận.
